% Source: https://artofproblemsolving.com/community/c3h3165919_2019_fall_lmt_individual_round__lexington_mathematical_tournament
% Stable identity: aops-topic-3165919-q13
Use $A=(0,0)$, $B=(2,0)$, $C=(2,2)$, $D=(0,2)$. Then $E=(2,1)$ and $F=(1,2)$, while the perpendicular line is $x+y=4$. Its intersections with $y=\dfrac{x}{2}$ and $y=2x$ are $X=(\dfrac{8}{3},\dfrac{4}{3})$ and $Y=(\dfrac{4}{3},\dfrac{8}{3})$. The determinant area is $\dfrac{1}{2}(\dfrac{64}{9}-\dfrac{16}{9})=\boxed{\dfrac{8}{3}}$.
\begin{asy}
usepackage("amsmath");
size(7cm);
pair A=(0,0),Bp=(2,0),C=(2,2),D=(0,2),Ep=(2,1),F=(1,2),X=(8/3,4/3),Y=(4/3,8/3);draw(A--Bp--C--D--cycle,gray);draw(A--X--Y--cycle,blue);label("$A$",A,SW);label("$B$",Bp,SE);label("$C$",C,NE);label("$D$",D,NW);label("$E$",Ep,E);label("$F$",F,W);label("$X$",X,E);label("$Y$",Y,N);
\end{asy}
